\begin{definition}[Global two- and three-site sector coordinates]
    \label{def:mpdo_global_sector_closure_coordinates}
    \lean{MPOTensor.PhysicalSectorFactorization.twoSiteGlobalRegroupEquiv,
      MPOTensor.PhysicalSectorFactorization.twoSiteGlobalRegroupEquiv_symm_apply,
      MPOTensor.PhysicalSectorFactorization.threeSiteGlobalRegroupEquiv,
      MPOTensor.PhysicalSectorFactorization.threeSiteGlobalRegroupEquiv_symm_apply}
    \leanok
    \uses{def:mpdo_two_site_sector_closure,
      def:mpdo_three_site_sector_closure}
    The two- and three-site physical spaces are reindexed by their canonical
    direct sums over sector pairs and triples, with each summand further
    regrouped into its boundary and neighboring factors. Write
    $\mathfrak G_2$ and $\mathfrak G_3$ for these reindexings of matrices.
\end{definition}

\begin{theorem}[Global direct-sum form of the sector closures]
    \label{thm:mpdo_global_sector_closure_decomposition}
    \lean{MPOTensor.PhysicalSectorFactorization.twoSiteGlobalClosure_eq,
      MPOTensor.PhysicalSectorFactorization.threeSiteGlobalClosure_eq}
    \leanok
    \uses{thm:mpdo_two_site_sector_closure_factorization,
      thm:mpdo_three_site_sector_closure_factorization,
      def:mpdo_physical_sector_factorization,
      def:mpdo_global_sector_closure_coordinates}
    After the canonical regrouping of all physical-sector indices, the full
    closures are
    \begin{align}
      \mathfrak G_2(\mc{K}_2(X))
      &=\bigoplus_{k,h}B_{k,h}(X)\otimes\eta_{k,h},
      \notag\\
      \mathfrak G_3(\mc{K}_3(X))
      &=\bigoplus_{k,l,h}B_{k,h}(X)\otimes
        (\eta_{k,l}\otimes\eta_{l,h}).
      \notag
    \end{align}
    Every matrix entry between distinct sector tuples vanishes.
\end{theorem}

\begin{proof}
    \leanok
    \uses{thm:mpdo_two_site_sector_closure_factorization,
      thm:mpdo_three_site_sector_closure_factorization,
      def:mpdo_physical_sector_factorization,
      def:mpdo_global_sector_closure_coordinates}
    On a diagonal sector tuple these are the fixed-sector factorizations.
    For distinct sector labels $k\ne p$, the block structure of the
    sector-coordinate tensor gives
    $[U\kappa_{\beta,\alpha}U^\dagger]_{k,p}=0$.
    Consequently, if $(k,h)\ne(p,q)$ or
    $(k,l,h)\ne(p,m,q)$, respectively, then
    \begin{align}
      [\mathfrak G_2(\mc{K}_2(X))]_{(k,h),(p,q)}
      &=0,
      \notag\\
      [\mathfrak G_3(\mc{K}_3(X))]_{(k,l,h),(p,m,q)}
      &=0.
      \notag
    \end{align}
\end{proof}

\begin{definition}[Two-site sector bond]
    \label{def:mpdo_physical_sector_two_site_bond}
    \lean{MPOTensor.PhysicalSectorFactorization.sectorCoordinateBond}
    \leanok
    \uses{def:mpdo_minimal_neighboring_operator,
      def:mpdo_global_sector_closure_coordinates}
    In the two-site sector coordinates, define
    \begin{align}
      B_{\mathrm{sector}}
      &=\bigoplus_{k,h}
        \left(\Id_{B_k^L\otimes B_h^R}\otimes\eta_{k,h}\right).
      \notag
    \end{align}
    Thus the outer factors carry the identity, while the neighboring factors
    carry the operator $\eta_{k,h}$. This is the local bond in
    \cite[Appendix~C.2, Proposition~C.8, lines~1581--1593]
      {Cirac2016MPDO_arXiv}.
\end{definition}

\begin{theorem}[Positivity of the two-site sector bond]
    \label{thm:mpdo_physical_sector_two_site_bond_pos}
    \lean{MPOTensor.PhysicalSectorFactorization.sectorCoordinateBond_pos}
    \leanok
    \uses{def:mpdo_physical_sector_two_site_bond,
      thm:matrix_pos_semidef_block_diagonal}
    Suppose that $\eta_{k,h}$ is positive semidefinite for every pair of
    sectors. Then $B_{\mathrm{sector}}$ is positive semidefinite. The positivity
    premise is the conclusion at
    \cite[Appendix~C.2, lines~1446--1450]{Cirac2016MPDO_arXiv}.
\end{theorem}

\begin{proof}
    \leanok
    \uses{thm:matrix_pos_semidef_block_diagonal}
    Each summand
    $\Id_{B_k^L\otimes B_h^R}\otimes\eta_{k,h}$ is positive semidefinite.
    Positivity is preserved by finite block diagonals and by the reindexing
    from the regrouped direct sum to the two-site sector coordinates.
\end{proof}

\begin{definition}[Two-site bond in physical coordinates]
    \label{def:mpdo_physical_two_site_bond}
    \lean{MPOTensor.PhysicalSectorFactorization.physicalPairBond,
      MPOTensor.PhysicalSectorFactorization.physicalBond}
    \leanok
    \uses{def:mpdo_physical_sector_two_site_bond,
      def:mpdo_physical_sector_factorization}
    Let $V_2$ be the two-site tensor product of the project's coordinate map,
    so that $V_2XV_2^\dagger$ expresses a physical operator $X$ in sector
    coordinates. Thus $V_2=U_{\mathrm{paper}}^\dagger$ in the convention of
    \cite[Appendix~C.2, lines~1581--1583]{Cirac2016MPDO_arXiv}. Define
    \begin{align}
      B_{\mathrm{physical}}
      &=V_2^\dagger B_{\mathrm{sector}}V_2,
      \notag
    \end{align}
    and identify pairs of physical indices with functions
    $\{0,1\}\to\{0,\ldots,d-1\}$. This is the physical two-site bond in
    \cite[Appendix~C.2, Proposition~C.8, lines~1581--1593]
      {Cirac2016MPDO_arXiv}.
\end{definition}

\begin{theorem}[Positivity of the physical two-site bond]
    \label{thm:mpdo_physical_two_site_bond_pos}
    \lean{MPOTensor.PhysicalSectorFactorization.physicalPairBond_pos,
      MPOTensor.PhysicalSectorFactorization.physicalBond_pos}
    \leanok
    \uses{def:mpdo_physical_two_site_bond,
      thm:mpdo_physical_sector_two_site_bond_pos}
    Suppose that $\eta_{k,h}$ is positive semidefinite for every pair of
    sectors. Then $B_{\mathrm{physical}}$ is positive semidefinite, both with
    pair indices and with two-site configuration indices.
\end{theorem}

\begin{proof}
    \leanok
    \uses{thm:mpdo_physical_sector_two_site_bond_pos}
    Apply positivity of $B_{\mathrm{sector}}$ to the congruence
    $V_2^\dagger B_{\mathrm{sector}}V_2$. Reindexing the resulting matrix by
    two-site configurations preserves positivity.
\end{proof}

\begin{definition}[Three-site lifts of a two-site operator]
    \label{def:mpdo_three_site_pair_lifts}
    \lean{MPOTensor.PhysicalSectorFactorization.leftPairMatrix,
      MPOTensor.PhysicalSectorFactorization.rightPairMatrix}
    \leanok
    Let $B$ be an operator on two copies of a finite-dimensional space $H$.
    On the right-associated space $H\otimes(H\otimes H)$, define
    \begin{align}
      B_{01}
      &=\alpha(B\otimes\Id_H)\alpha^{-1},
      &
      B_{12}
      &=\Id_H\otimes B,
      \notag
    \end{align}
    where $\alpha:(H\otimes H)\otimes H\to H\otimes(H\otimes H)$ is the
    canonical associator.
\end{definition}

\begin{theorem}[Three-site translation identities]
    \label{thm:mpdo_three_site_pair_lifts}
    \lean{MPOTensor.PhysicalSectorFactorization.reindex_embedLocalOperator_zero,
      MPOTensor.PhysicalSectorFactorization.reindex_embedLocalOperator_one}
    \leanok
    \uses{def:mpdo_physical_two_site_bond,
      def:mpdo_three_site_pair_lifts}
    Under the canonical identification of three-site configurations with
    $H\otimes(H\otimes H)$, the translates of the physical bond beginning at
    sites $0$ and $1$ are respectively
    $(B_{\mathrm{physical}})_{01}=B_{01}$ and
    $(B_{\mathrm{physical}})_{12}=B_{12}$.
\end{theorem}

\begin{proof}
    \leanok
    An entry of the first translate vanishes unless the third indices agree;
    its remaining entry is that of $B_{\mathrm{physical}}$. Similarly, an
    entry of the second translate vanishes unless the first indices agree.
    These are precisely the matrix entries of $B_{01}$ and $B_{12}$.
\end{proof}

\begin{theorem}[Commutativity on a fixed sector triple]
    \label{thm:mpdo_fixed_sector_adjacent_bonds_comm}
    \lean{MPOTensor.PhysicalSectorFactorization.adjacentSectorBonds_comm}
    \leanok
    \uses{def:mpdo_minimal_neighboring_operator}
    Fix sectors $(k,l,h)$. On the regrouped sector space, set
    \begin{align}
      C_{01}^{k,l,h}
      &=\Id_{B_k^L\otimes B_h^R}\otimes
        \left(\eta_{k,l}\otimes
        \Id_{B_l^R\otimes B_h^L}\right),
      \notag\\
      C_{12}^{k,l,h}
      &=\Id_{B_k^L\otimes B_h^R}\otimes
        \left(\Id_{B_k^R\otimes B_l^L}\otimes\eta_{l,h}\right).
      \notag
    \end{align}
    Then
    $C_{01}^{k,l,h}C_{12}^{k,l,h}
    =C_{12}^{k,l,h}C_{01}^{k,l,h}$.
    This is the fixed-sector calculation in
    \cite[Appendix~C.2, Proposition~C.8, lines~1589--1593]
      {Cirac2016MPDO_arXiv}.
\end{theorem}

On the sector triple $(k,l,h)$, the two operators act on distinct neighboring
factors, while the outer factors $B_k^L$ and $B_h^R$ remain unchanged. Their
commutativity is therefore the local identity
$\eta_{k,l}^{(12)}\eta_{l,h}^{(23)}
=\eta_{l,h}^{(23)}\eta_{k,l}^{(12)}$.

\begin{proof}
    \leanok
    The operators $\eta_{k,l}$ and $\eta_{l,h}$ act on the distinct factors
    $B_k^R\otimes B_l^L$ and $B_l^R\otimes B_h^L$, respectively. The claim
    follows from the mixed-product identity for tensor products.
\end{proof}

\begin{theorem}[Commutativity of adjacent physical bonds]
    \label{thm:mpdo_physical_adjacent_bonds_comm}
    \lean{MPOTensor.PhysicalSectorFactorization.physicalBond_zero_one_comm}
    \leanok
    \uses{def:mpdo_physical_two_site_bond,
      def:mpdo_global_sector_closure_coordinates,
      thm:mpdo_three_site_pair_lifts,
      thm:mpdo_fixed_sector_adjacent_bonds_comm}
    The two translates of the physical bond on three sites commute:
    \begin{align}
      (B_{\mathrm{physical}})_{01}(B_{\mathrm{physical}})_{12}
      &=(B_{\mathrm{physical}})_{12}(B_{\mathrm{physical}})_{01}.
      \notag
    \end{align}
    This is the adjacent-bond calculation in
    \cite[Appendix~C.2, Proposition~C.8, lines~1589--1593]
      {Cirac2016MPDO_arXiv}.
\end{theorem}

\begin{proof}
    \leanok
    Regroup the three-site sector coordinates into the direct sum over triples
    $(k,l,h)$. The two bonds become block diagonal, with blocks
    $C_{01}^{k,l,h}$ and $C_{12}^{k,l,h}$, which commute by
    Theorem~\ref{thm:mpdo_fixed_sector_adjacent_bonds_comm}. Conjugating by the
    three-fold physical coordinate unitary preserves the equality. The
    three-site translation identities then give the displayed physical
    commutator.
\end{proof}

\begin{theorem}[Agreement outside a cyclic window]
    \label{thm:mpdo_agreement_outside_cyclic_window}
    \lean{MPOTensor.agreesOutsideWindow_iff}
    \leanok
    Let $W_i^L=\{i,i+1,\ldots,i+L-1\}$ be a cyclic window of length
    $L\leq N$. Two configurations agree outside $W_i^L$ precisely when their
    values agree at every site not contained in $W_i^L$.
\end{theorem}

\begin{proof}\leanok
    Replacing the entries in $W_i^L$ does not affect any complementary site.
    Conversely, if the two configurations agree on the complement, replacing
    the window of the first by the window of the second reconstructs the
    second configuration.
\end{proof}

\begin{theorem}[Multiplicativity of cyclic local embeddings]
    \label{thm:mpdo_cyclic_local_embedding_mul}
    \lean{MPOTensor.embedLocalOperator_mul}
    \leanok
    Let $E_i^L$ embed an operator on $L$ consecutive sites into an
    $N$-site periodic chain at the cyclic window beginning at $i$. Then
    $E_i^L(BC)=E_i^L(B)E_i^L(C)$.
\end{theorem}

\begin{proof}\leanok
    Identify a periodic configuration with its restriction to $W_i^L$ and
    its restriction to the cyclic complement. In these coordinates,
    $E_i^L(B)=B\otimes\Id$, so the assertion is the mixed-product identity.
\end{proof}

\begin{theorem}[Commutativity on disjoint cyclic windows]
    \label{thm:mpdo_disjoint_cyclic_local_embeddings_comm}
    \lean{MPOTensor.embedLocalOperator_commute_of_not_cyclicWindowsOverlap}
    \leanok
    \uses{lem:cyclic_window_support_offsets,
      lem:cyclic_windows_overlap_symmetry}
    Let $L\leq N$, and let $W_i^L$ and $W_j^L$ be disjoint cyclic windows in
    an $N$-site periodic chain. For arbitrary operators $B$ and $C$ on $L$
    sites, $E_i^L(B)E_j^L(C)=E_j^L(C)E_i^L(B)$.
\end{theorem}

\begin{proof}\leanok
    Apply both sides to a chain amplitude. The first order of application is
    a double sum over replacements on $W_i^L$ and $W_j^L$. Since the windows
    are disjoint, each replacement leaves the other restricted configuration
    unchanged, and the two replacements commute. Interchanging the two finite
    sums gives the reverse order.
\end{proof}

\begin{theorem}[Adjacent physical bonds on periodic chains]
    \label{thm:mpdo_periodic_adjacent_physical_bonds_comm}
    \lean{MPOTensor.PhysicalSectorFactorization.physicalBond_adjacent_comm}
    \leanok
    \uses{thm:mpdo_cyclic_local_embedding_mul,
      thm:mpdo_physical_adjacent_bonds_comm}
    Let $N\geq 3$. For every site $i\in\Z/N\Z$, including
    $i=N-1$, the adjacent translates of the physical bond commute:
    \begin{align}
      (B_{\mathrm{physical}})_{i,i+1}
      (B_{\mathrm{physical}})_{i+1,i+2}
      &=
      (B_{\mathrm{physical}})_{i+1,i+2}
      (B_{\mathrm{physical}})_{i,i+1}.
      \notag
    \end{align}
    This is the periodic-chain form of the adjacent-bond calculation in
    \cite[Appendix~C.2, Proposition~C.8, lines~1571--1593]
      {Cirac2016MPDO_arXiv}. The case $N=2$ is treated separately below,
    because its two translated bonds have the same support with opposite
    cyclic order.
\end{theorem}

\begin{proof}
    \leanok
    Enclose the two bonds in the cyclic three-site window beginning at $i$.
    In these coordinates they are the bonds beginning at sites $0$ and $1$,
    respectively. Embedding an operator into a fixed cyclic window preserves
    products, so the assertion follows from
    Theorem~\ref{thm:mpdo_physical_adjacent_bonds_comm}.
\end{proof}

\begin{theorem}[Commutativity of the crossed two-site translates]
    \label{thm:mpdo_physical_two_site_crossed_bonds_comm}
    \lean{MPOTensor.PhysicalSectorFactorization.physicalBond_two_zero_one_comm}
    \leanok
    \uses{def:mpdo_global_sector_closure_coordinates,
      def:mpdo_physical_sector_two_site_bond,
      def:mpdo_physical_two_site_bond}
    On the periodic chain of length two, the translates beginning at sites
    zero and one read the physical sites in the opposite cyclic orders
    $(0,1)$ and $(1,0)$. Nevertheless they commute:
    $B_{01}B_{10}=B_{10}B_{01}$.
    This is the length-two instance of the translated-bond commutativity in
    \cite[Appendix~C.2, Proposition~C.8]{Cirac2016MPDO_arXiv}. The source
    does not discuss the crossed finite-size ordering separately.
\end{theorem}

\begin{proof}
    \leanok
    Fix sectors $(k,h)$ and regroup the two physical sites as
    $(L_k\otimes R_h)\otimes(R_k\otimes L_h)$.
    The translate in the order $(0,1)$ is
    $\Id_{L_k\otimes R_h}\otimes\eta_{k,h}$, while the translate in the
    order $(1,0)$ is
    $\eta_{h,k}^{\mathrm{op}}\otimes\Id_{R_k\otimes L_h}$, where
    $\eta_{h,k}^{\mathrm{op}}$ denotes the same matrix after exchanging
    the factors $R_h$ and $L_k$. These operators act on complementary
    factors and commute. Taking the direct sum over $(k,h)$ and conjugating
    by the two-fold physical coordinate unitary proves the claim.
\end{proof}

\begin{theorem}[Disjoint physical bonds on periodic chains]
    \label{thm:mpdo_periodic_disjoint_physical_bonds_comm}
    \lean{MPOTensor.PhysicalSectorFactorization.physicalBond_disjoint_comm}
    \leanok
    \uses{def:mpdo_physical_two_site_bond,
      thm:mpdo_disjoint_cyclic_local_embeddings_comm}
    Let $N\geq 2$, and suppose that the cyclic bonds beginning at $i$ and $j$
    have disjoint two-site supports. Then their physical bond operators
    commute:
    $(B_{\mathrm{physical}})_{i,i+1}
    (B_{\mathrm{physical}})_{j,j+1}
    =(B_{\mathrm{physical}})_{j,j+1}
    (B_{\mathrm{physical}})_{i,i+1}$.
    This is the locality part of the pairwise commutation assertion in
    \cite[Appendix~C.2, Proposition~C.8, lines~1571--1593]
      {Cirac2016MPDO_arXiv}.
\end{theorem}

\begin{proof}\leanok
    This is Theorem~\ref{thm:mpdo_disjoint_cyclic_local_embeddings_comm}
    applied to the same physical bond operator in both windows. No positivity
    or additional property of the sector factorization is needed for this
    locality step.
\end{proof}

\begin{theorem}[Pairwise commutativity of the physical bonds]
    \label{thm:mpdo_periodic_pairwise_physical_bonds_comm}
    \lean{MPOTensor.PhysicalSectorFactorization.physicalBond_pairwise_comm}
    \leanok
    \uses{lem:two_site_cyclic_window_overlap_cases,
      thm:mpdo_periodic_adjacent_physical_bonds_comm,
      thm:mpdo_physical_two_site_crossed_bonds_comm,
      thm:mpdo_periodic_disjoint_physical_bonds_comm}
    Suppose that the physical-sector factorization is given. For every
    periodic chain of length $N\geq2$ and every pair of sites
    $i,j\in\Z/N\Z$, the corresponding translates of the
    physical bond commute:
    \begin{align}
      (B_{\mathrm{physical}})_{i,i+1}
      (B_{\mathrm{physical}})_{j,j+1}
      &=
      (B_{\mathrm{physical}})_{j,j+1}
      (B_{\mathrm{physical}})_{i,i+1}.
      \notag
    \end{align}
    This is the pairwise commutation assertion in
    \cite[Appendix~C.2, Proposition~C.8, lines~1571--1593]
      {Cirac2016MPDO_arXiv}, conditional here on the physical-sector
    factorization constructed in the preceding results.
\end{theorem}

\begin{proof}
    \leanok
    For $N=2$, equal translates commute trivially, while the two distinct
    translates commute by
    Theorem~\ref{thm:mpdo_physical_two_site_crossed_bonds_comm}. Let $N\geq3$.
    If the two cyclic supports are disjoint, apply
    Theorem~\ref{thm:mpdo_periodic_disjoint_physical_bonds_comm}. Otherwise,
    Lemma~\ref{lem:two_site_cyclic_window_overlap_cases} shows that the two
    starting sites either coincide or are cyclic neighbors. The coincident
    case is immediate, and the two neighboring orientations follow from
    Theorem~\ref{thm:mpdo_periodic_adjacent_physical_bonds_comm} and its
    symmetric equality.
\end{proof}

\begin{theorem}[Exact product of physical-sector bonds]
    \label{thm:mpdo_exact_physical_sector_bond_product}
    \lean{MPOTensor.PhysicalSectorFactorization.mpo_eq_product_physicalBond}
    \leanok
    \uses{def:mpdo_physical_sector_factorization,
      def:mpdo_physical_two_site_bond,
      thm:mpdo_physical_sector_chain_decomposition}
    Suppose that $K$ is equipped with a physical-sector factorization $F$.
    Let $B_i$ denote the cyclic translate beginning at site $i$ of the
    physical two-site bond determined by $F$. For every $N\geq2$, the
    $N$-site periodic operator is
    \begin{align}
      \rho_N
      &=B_0B_1\cdots B_{N-1}.
      \notag
    \end{align}
    No positivity, zero-correlation-length, injectivity, or normalization
    hypothesis is required. This is the product identity in
    \cite[Appendix~C.2, Proposition~C.8, lines~1581--1593]
      {Cirac2016MPDO_arXiv}, conditional here on the given physical-sector
    factorization.
\end{theorem}

\begin{proof}
    \leanok
    For a fixed cyclic sector sequence, contraction of the virtual index
    between the right tensor at site $i$ and the left tensor at site $i+1$
    produces $\eta_{k_i,k_{i+1}}$. Taking the virtual trace therefore gives
    the cyclic product of these neighboring operators. The
    translated sector bonds give the same ordered product. Conjugating by the
    tensor power of the one-site coordinate unitary gives the stated physical
    identity. For $N=2$, the second translate has the opposite cyclic order;
    the two factors are the ordinary and crossed bonds.
\end{proof}

\begin{theorem}[Unit scalar in the physical-sector product]
    \label{thm:mpdo_physical_sector_bond_product_unit_scalar}
    \lean{MPOTensor.PhysicalSectorFactorization.exists_positive_scalar_mpo_eq_product_physicalBond}
    \leanok
    \uses{thm:mpdo_exact_physical_sector_bond_product}
    Under the physical-sector factorization $F$ of the preceding theorem,
    for every $N\geq2$ there is a real number $c>0$ such that
    $\rho_N=c\,B_0B_1\cdots B_{N-1}$.
    Here the bonds are the canonical bonds determined by $F$, and one may
    take $c=1$. No positivity hypothesis on the neighboring operators
    $\eta_{k,h}$ is required.
\end{theorem}

\begin{proof}
    \leanok
    Take $c=1$ in
    Theorem~\ref{thm:mpdo_exact_physical_sector_bond_product}.
\end{proof}

\begin{theorem}[Positive translation-invariant bond datum]
    \label{thm:mpdo_positive_translation_invariant_bond_data}
    \lean{MPOTensor.PhysicalSectorFactorization.translationInvariantBondData}
    \leanok
    \uses{thm:mpdo_physical_two_site_bond_pos,
      thm:mpdo_periodic_pairwise_physical_bonds_comm}
    Suppose that every neighboring operator $\eta_{k,h}$ is positive
    semidefinite. Then the physical bond is positive semidefinite and all of
    its cyclic translates commute on every periodic chain of length at least
    two. Thus it determines a single translation-invariant positive bond
    datum. This assertion does not include the product formula for the
    finite-chain density operator.
\end{theorem}

\begin{proof}
    \leanok
    Positivity is
    Theorem~\ref{thm:mpdo_physical_two_site_bond_pos}; pairwise commutativity
    is Theorem~\ref{thm:mpdo_periodic_pairwise_physical_bonds_comm}.
\end{proof}

\begin{theorem}[$\eta$-local structure from a positive physical-sector factorization]
    \label{thm:mpdo_eta_local_structure_of_positive_physical_sector_factorization}
    \lean{MPOTensor.PhysicalSectorFactorization.etaLocalStructureData}
    \lean{MPOTensor.PhysicalSectorFactorization.fixedProductTensorData}
    \leanok
    \uses{thm:mpdo_exact_physical_sector_bond_product,
      thm:mpdo_positive_translation_invariant_bond_data,
      def:mpdo_eta_local_structure_data}
    Suppose that a physical-sector factorization of $K$ is given and that
    every neighboring operator $\eta_{k,h}$ is positive semidefinite. Then
    the physical two-site bond determines an $\eta$-local structure for $K$.
    Its translated copies commute pairwise, and for every $N\geq2$,
    $\rho_N=B_0B_1\cdots B_{N-1}$.
    If the virtual dimension of $K$ is positive, then $K$ itself is a fixed
    tensor representation of this product family, with the same virtual
    dimension. Thus the positive normalization scalar may be chosen to be
    $1$. This is Proposition~C.8 conditional on the coherently positive
    physical-sector factorization that the source derives from SAL.
\end{theorem}

\begin{proof}
    \leanok
    Theorem~\ref{thm:mpdo_positive_translation_invariant_bond_data} gives the
    positive translation-invariant bond and its pairwise commuting
    translates. Theorem~\ref{thm:mpdo_exact_physical_sector_bond_product}
    realizes every finite-chain operator with normalization scalar $1$.
\end{proof}

\begin{corollary}[$\eta$-local structure from SAL]%
    \label{cor:mpdo_eta_local_structure_of_is_sal}
    \lean{MPOTensor.nonempty_etaLocalStructureData_of_isSAL}
    \leanok
    \uses{thm:mpdo_eta_local_structure_of_positive_physical_sector_factorization,
      thm:mpdo_eta_coherent_positive_choice}
    Every injective MPO tensor satisfying SAL admits an $\eta$-local
    structure. Thus there is a single positive two-site bond whose cyclic
    translates commute pairwise and whose ordered product is the finite-chain
    density operator at every length $N\geq2$, with normalization scalar one.
    This is the conclusion of Proposition~C.8.
\end{corollary}

\begin{proof}
    \leanok
    Theorem~\ref{thm:mpdo_eta_coherent_positive_choice} supplies the positive
    physical-sector factorization. Apply
    Theorem~\ref{thm:mpdo_eta_local_structure_of_positive_physical_sector_factorization}.
\end{proof}

\begin{corollary}[Normal representative for the SAL-selected bond]%
    \label{cor:mpdo_normal_fixed_product_tensor_of_is_sal}
    \lean{MPOTensor.exists_normal_fixedProductTensorData_of_isSAL}
    \leanok
    \uses{def:mpdo_fixed_bond_product_tensor,
      def:mpdo_physical_sector_factorization, def:nt_cpsv}
    Let $K$ be injective and satisfy SAL, and suppose that its doubled-index
    matrix-product tensor is normal. There is a positive physical-sector
    factorization whose selected two-site bond $B$ has $K$ itself as a normal
    fixed tensor representation. Its virtual dimension is the virtual
    dimension of $K$, and for every $N\geq2$,
    $\rho_N(K)=B_0B_1\cdots B_{N-1}$.
    In particular, the scalar is one. This is an existential statement about
    the bond selected in the proof of Proposition~C.8; it does not assert
    uniqueness or normal rescalability for an arbitrary proportional
    commuting-bond presentation.
\end{corollary}

\begin{proof}
    \leanok
    \uses{thm:mpdo_exact_physical_sector_bond_product,
      thm:mpdo_eta_coherent_positive_choice}
    SAL implies that the virtual dimension is positive. The positive
    physical-sector factorization is supplied by
    Theorem~\ref{thm:mpdo_eta_coherent_positive_choice}. For its selected
    bond, use $K$ as the fixed tensor. The exact product identity is
    Theorem~\ref{thm:mpdo_exact_physical_sector_bond_product}, and normality
    is the given hypothesis on the doubled-index tensor.
\end{proof}
